
\documentclass[11pt,a4paper]{article}
\usepackage[margin=2.2cm]{geometry}
\usepackage[utf8]{inputenc}
\usepackage[T1]{fontenc}
\usepackage{amsmath,amssymb,mathtools,bm}
\usepackage{graphicx,booktabs,longtable,array,microtype,xcolor,hyperref,fancyhdr,enumitem}
\graphicspath{{../images/}}
\definecolor{navy}{HTML}{17324D}\definecolor{rust}{HTML}{A33A2B}\definecolor{green}{HTML}{176B4D}
\hypersetup{colorlinks=true,linkcolor=navy,urlcolor=rust,citecolor=green}
\pagestyle{fancy}\fancyhf{}\lhead{Auditoria estrutural da RH}\rhead{Dossiê técnico}\cfoot{\thepage}
\newcommand{\Drad}{\mathfrak D_{\mathrm{rad}}}
\newcommand{\Dpar}{\Delta_{\mathrm{par}}}
\newcommand{\status}[2]{\par\noindent\colorbox{#1!12}{\textcolor{#1}{\textbf{#2}}}\quad}
\begin{document}

\begin{titlepage}\centering{\Huge\sffamily\bfseries Caderno de fórmulas e verificações\par}\vfill{\Large Detector radial e obrigações de prova\par}\vfill Agosto de 2026\end{titlepage}
\section{Parâmetros fundamentais}
\begin{align*}
\rho&=\tfrac12+\delta+i\gamma,& \rho^\sharp&=1-\bar\rho=\tfrac12-\delta+i\gamma,\\
R_c^2&=\gamma^2+\tfrac14,& R_\pm^2&=R_c^2\pm\delta+\delta^2,\\
A&=\gamma^2+\tfrac14+\delta^2,& P&=A^2-\delta^2.
\end{align*}
\section{Reconstrução exata}
\[\delta=\frac{R_+^2-R_-^2}{2},\qquad \delta^2=\frac{R_+^2+R_-^2}{2}-R_c^2.\]
\section{Positividade local}
\[w=\frac{2\delta^2}{P}\ge0,\qquad w=0\Longleftrightarrow\delta=0,\qquad W_4=4w.\]
Para $|\delta|<1/2$ e $\gamma\ne0$, o denominador é estritamente positivo porque é o produto $|\rho|^2|1-\rho|^2$.
\section{Expansão assintótica}
Pondo $a=\tfrac14+\delta^2$,
\begin{align*}
P&=\gamma^4+2a\gamma^2+(a^2-\delta^2),\\
P^{-1}&=\gamma^{-4}\left(1-2a\gamma^{-2}+O(\gamma^{-4})\right),\\
w&=2\delta^2\gamma^{-4}-4\delta^2a\gamma^{-6}+O(\gamma^{-8}).
\end{align*}
\section{Constante radial}
\[C_\xi=2+\gamma_E-\log(4\pi)\approx0.0461914179322420676.\]
\section{Defeito global}
\[\Drad=\sum_\rho\frac{2(\Re\rho-\tfrac12)^2}{|\rho|^2|1-\rho|^2}.\]
Cada parcela é não negativa. A dificuldade não está nessa positividade, mas em provar a anulação da soma por uma identidade independente que alcance o termo não holomorfo.
\section{Forma de Weil}
\[Q_W(g)=\sum_\rho\widetilde g(\rho)\overline{\widetilde g(1-\bar\rho)}.\]
A obrigação construtiva é produzir $T$ a partir do lado dos primos e provar $Q_W(g)=\|Tg\|^2$ sem assumir a positividade que se deseja concluir.
\section{Matrizes de Gram}
\[(G_N)_{ij}=Q_W(\phi_i,\phi_j),\qquad G_N=G_N^*,\qquad G_N\succeq0.\]
Uma conclusão global exige: (i) positividade em todos os níveis; (ii) união densa dos espaços; (iii) módulo de continuidade para $Q_W$; (iv) erros intervalares uniformes.
\section{Extensão por $\epsilon$}
\[\mathbb{RC}=\mathcal A[\epsilon]/(\epsilon^2-\Dpar).\]
Em geral $\epsilon$ não é nilpotente quando $\Dpar\ne0$. Declarar $\epsilon^2=0$ universalmente é equivalente a impor o setor crítico.
\section{Checklist de não circularidade}
\begin{enumerate}[label=\arabic*.]
\item Definir todos os objetos antes de assumir RH.
\item Repetir a construção num cenário com um zero sintético fora da linha.
\item Identificar exatamente qual lema falha nesse cenário.
\item Exigir que esse lema seja derivado de dados primos e arquimedianos.
\item Não converter certificados finitos em quantificadores universais.
\end{enumerate}

\end{document}
